\documentclass[JDR,PDF]{jdr} 
\usepackage{fancyhdr}
\usepackage{layout}
\usepackage{multirow}
\usepackage{makecell}
\usepackage{afterpage}
\usepackage{siunitx}
\usepackage{enumerate}
\usepackage{amsthm}
\usepackage{draftwatermark}
\SetWatermarkText{Confidential: For Review Only}
\SetWatermarkScale{2}
\SetWatermarkAngle{-45}
\newcommand*{\doi}[1]{\href{https://doi.org/\detokenize{#1}}{doi: \detokenize{#1}}}
\newenvironment{myproof}[2] {\paragraph{Proof of {#1} {#2}:}}{\hfill}

\begin{document}

\begin{picture}(0,0)
\put(400,8){\includegraphics[width=2.5cm]{JDRLogo.png}}
\end{picture}


\title{Title}

\articletype{Article Type} % Original Article, Review Article, Communication, Erratum


\author{
First Author\inst{1,2}\email{login1@email.ext},
Other Authors\inst{2}\email{login2@email.ext}
}

\institute{
\inst{1}
          First Author Affiliation, Address 
\inst{2}
          Other Authors Affiliation, Address 
}


\abstract{A single paragraph of about 200 words maximum. For research articles, abstracts should give a pertinent overview of the work. We strongly encourage authors to use the following style of structured abstracts, but without headings: (1) Background: Place the question addressed in a broad context and highlight the purpose of the study; (2) Methods: Describe briefly the main methods or treatments applied; (3) Results: Summarize the article's main findings; and (4) Conclusion: Indicate the main conclusions or interpretations. The abstract should be an objective representation of the article, it must not contain results which are not presented and substantiated in the main text and should not exaggerate the main conclusions.}

\keywords{keyword 1 \*\ keyword 2 \*\ keyword 3 (list three to ten pertinent keywords specific to the article, yet reasonably common within the subject discipline.)}


\setcounter{page}{3}
\setlength{\parindent}{20pt}
\maketitle

%%%%%%%%%%%%%%%%%%%%%%%%%%%%%%%%%%%%%%
%%%%%%%%%%%%Introduction%%%%%%%%%%%%%%
%%%%%%%%%%%%%%%%%%%%%%%%%%%%%%%%%%%%%%
\section{Introduction}

The introduction should briefly place the study in a broad context and highlight why it is important. It should define the purpose of the work and its significance. The current state of the research field should be reviewed carefully and key publications cited. Please highlight controversial and diverging hypotheses when necessary. Finally, briefly mention the main aim of the work and highlight the principal conclusions. As far as possible, please keep the introduction comprehensible to scientists outside your particular field of research. Citing a journal paper \cite{ref-journal}. And now citing a book reference \cite{ref-book}. 

%%%%%%%%%%%%%%%%%%%%%%%%%%%%%%%%%%%%%%
%%%%%%%%%%%%%%Methods%%%%%%%%%%%%%%%%%
%%%%%%%%%%%%%%%%%%%%%%%%%%%%%%%%%%%%%%
\section{Methods}

Materials and Methods should be described with sufficient details to allow others to replicate and build on published results. Please note that publication of your manuscript implicates that you must make all materials, data, computer code, and protocols associated with the publication available to readers. Please disclose at the submission stage any restrictions on the availability of materials or information. New methods and protocols should be described in detail while well-established methods can be briefly described and appropriately cited.

Research manuscripts reporting large datasets that are deposited in a publicly available database should specify where the data have been deposited and provide the relevant accession numbers. If the accession numbers have not yet been obtained at the time of submission, please state that they will be provided during review. They must be provided prior to publication.

Interventionary studies involving animals or humans, and other studies require ethical approval must list the authority that provided approval and the corresponding ethical approval code. 
 
\begin{quote}
This section may be divided by subheadings. It should provide a concise and precise description of the experimental results, their interpretation as well as the experimental conclusions that can be drawn.
\end{quote}

%%%%%%%%%%%%%%%%%%%%%%%%%%%%%%%%%%%%%%%%%%
\subsection{Subsection}
A section may be further divided by subheadings.
\subsubsection{Subsubsection}

Bulleted lists look like this:
\begin{itemize}
\item First level, itemize, first item
\begin{itemize}
\item Second level, itemize, first item
\item Second level, itemize, second item
\begin{enumerate}
\item Third level, enumerate, first item
\item Third level, enumerate, second item
\end{enumerate}
\end{itemize}
\item First level, itemize, second item
\end{itemize}

Numbered lists can be added as follows:
\begin{enumerate}
\item	First item 
\item	Second item
\item	Third item
\end{enumerate}

The text continues here.

\subsection{Figures, Tables and Schemes}

All figures and tables should be cited in the main text as Fig.~\ref{fig:1}, Table~\ref{tab:1}, etc.

\begin{figure}[hbt!]
\centering
\includegraphics[width=6cm]{JDRLogo.png}
\caption{This is a figure, Schemes follow the same formatting. If there are multiple panels, they should be listed as: (\textbf{a}) Description of what is contained in the first panel. (\textbf{b}) Description of what is contained in the second panel. Figures should be placed in the main text near to the first time they are cited. A caption on a single line should be centered.}
\label{fig:1}
\end{figure}   
 

\begin{table}[hbt!]
\caption{This is a table caption. Tables should be placed in the main text near to the first time they are cited.}
\centering
%% \tablesize{} %% You can specify the fontsize here, e.g., \tablesize{\footnotesize}. If commented out \small will be used.
\begin{tabular}{ccc}
\hline
\textbf{Title 1}	& \textbf{Title 2}	& \textbf{Title 3}\\
\hline
entry 1		& data			& data\\
entry 2		& data			& data\\
\hline
\end{tabular}
\label{tab:1}
\end{table}

\subsection{Formatting of Mathematical Components}

This is an example of an equation:

\begin{equation} \label{eq:1}
\sum_{i=0}^{\infty} a x^i,
\end{equation}
where $a$ is a constant and $x$ is a variable.

Cross referencing equations is done as follows. The equation (\ref{eq:1}) is a typical power series.

Please punctuate equations as regular text. Theorem-type environments (including propositions, lemmas, corollaries etc.) can be formatted as follows:
%% Example of a theorem:
\begin{theorem} \label{thm:1}
Example text of a theorem. Example: Let $f$ be a function whose derivative exists in every point, then $f$ is a continuous function.
\end{theorem}

Similarly, cross referencing theorems goes as follows. The theorem (\ref{thm:1}) is a general proposition not self-evident but proved by a chain of reasoning.

The text continues here. Proofs must be formatted as follows:

%% Example of a proof:
\begin{myproof}{Theorem}{\ref{thm:1}}
Text of the proof.
\end{myproof}

The text continues here.




%%%%%%%%%%%%%%%%%%%%%%%%%%%%%%%%%%%%%%
%%%%%%%%%%%%%%Results%%%%%%%%%%%%%%%%%
%%%%%%%%%%%%%%%%%%%%%%%%%%%%%%%%%%%%%%
\section{Results}

This section may be divided by subheadings, too. It should provide a concise and precise description of the experimental results, their interpretation as well as the experimental conclusions that can be drawn.


%%%%%%%%%%%%%%%%%%%%%%%%%%%%%%%%%%%%%%
%%%%%%%%%%%%%Discussion%%%%%%%%%%%%%%%
%%%%%%%%%%%%%%%%%%%%%%%%%%%%%%%%%%%%%%
\section{Discussion}

Authors should discuss the results and how they can be interpreted in perspective of previous studies and of the working hypotheses. The findings and their implications should be discussed in the broadest context possible. Future research directions may also be highlighted.

%%%%%%%%%%%%%%%%%%%%%%%%%%%%%%%%%%%%%%
%%%%%%%%%%%%Conclusions%%%%%%%%%%%%%%%
%%%%%%%%%%%%%%%%%%%%%%%%%%%%%%%%%%%%%%
\section{Conclusions}

This section completes the discussion.

\section*{Conflict of Interest}\noindent
Authors of this article declare that they have no conflict of interest.

\section*{Human Studies/Informed Consent}\noindent
No human studies were carried out by the authors for this article.

\section*{Animal Studies}\noindent
No animal studies were carried out by the authors for this article.


\section*{References}
\begin{thebibliography}{9}
\vspace{-7mm}

% Reference 1
\bibitem[Author1(year)]{ref-journal}
Author1, T. The title of the cited article. {\em Journal Abbreviation} {\bf 2020}, {\em 10}, 142--149. \doi{digital.object.identifier}

% Reference 2
\bibitem[Author2(year)]{ref-book}
Author2, L. The title of the cited contribution. In {\em The Book Title}; Editor1, F., Editor2, A., Eds.; Publishing House: City, Country, 2007; pp. 32--58.

\end{thebibliography}

%\afterpage{\null\newpage}

\end{document}
